% Options for packages loaded elsewhere
\PassOptionsToPackage{unicode}{hyperref}
\PassOptionsToPackage{hyphens}{url}
\PassOptionsToPackage{dvipsnames,svgnames,x11names}{xcolor}
%
\documentclass[
  10pt,
]{article}
\usepackage{amsmath,amssymb}
\usepackage{iftex}
\ifPDFTeX
  \usepackage[T1]{fontenc}
  \usepackage[utf8]{inputenc}
  \usepackage{textcomp} % provide euro and other symbols
\else % if luatex or xetex
  \usepackage{unicode-math} % this also loads fontspec
  \defaultfontfeatures{Scale=MatchLowercase}
  \defaultfontfeatures[\rmfamily]{Ligatures=TeX,Scale=1}
\fi
\usepackage{lmodern}
\ifPDFTeX\else
  % xetex/luatex font selection
  \setmainfont[]{DejaVu Serif}
  \setsansfont[]{DejaVu Sans}
  \setmonofont[]{DejaVu Sans Mono}
  \setmathfont[]{Latin Modern Math}
\fi
% Use upquote if available, for straight quotes in verbatim environments
\IfFileExists{upquote.sty}{\usepackage{upquote}}{}
\IfFileExists{microtype.sty}{% use microtype if available
  \usepackage[]{microtype}
  \UseMicrotypeSet[protrusion]{basicmath} % disable protrusion for tt fonts
}{}
\makeatletter
\@ifundefined{KOMAClassName}{% if non-KOMA class
  \IfFileExists{parskip.sty}{%
    \usepackage{parskip}
  }{% else
    \setlength{\parindent}{0pt}
    \setlength{\parskip}{6pt plus 2pt minus 1pt}}
}{% if KOMA class
  \KOMAoptions{parskip=half}}
\makeatother
\usepackage{xcolor}
\usepackage[margin=2.0cm]{geometry}
\usepackage{longtable,booktabs,array}
\usepackage{calc} % for calculating minipage widths
% Correct order of tables after \paragraph or \subparagraph
\usepackage{etoolbox}
\makeatletter
\patchcmd\longtable{\par}{\if@noskipsec\mbox{}\fi\par}{}{}
\makeatother
% Allow footnotes in longtable head/foot
\IfFileExists{footnotehyper.sty}{\usepackage{footnotehyper}}{\usepackage{footnote}}
\makesavenoteenv{longtable}
\setlength{\emergencystretch}{3em} % prevent overfull lines
\providecommand{\tightlist}{%
  \setlength{\itemsep}{0pt}\setlength{\parskip}{0pt}}
\setcounter{secnumdepth}{-\maxdimen} % remove section numbering
% doc_header.tex -- LaTeX preamble additions used by build_docs.sh for every
% Markdown -> LaTeX -> PDF conversion (README and provenance documents of Revision_SuperLoop).
\usepackage{etoolbox}
% longtable may not be loaded by pandoc when a document has no table
\usepackage{longtable}
% bold vectors: \bm -> \symbf under unicode-math (XeLaTeX), bm package otherwise
\makeatletter
\@ifpackageloaded{unicode-math}{\providecommand{\bm}[1]{\symbf{#1}}}{\usepackage{bm}}
\makeatother
% smaller type in long tables so that wide audit tables fit the text width;
% inside a table a line break is allowed after an underscore of plain text
% (pandoc writes it as \_), so that long identifiers wrap inside their column
% rows are set a little apart, so that symbols with indices do not touch
\AtBeginEnvironment{longtable}{\footnotesize\renewcommand{\arraystretch}{1.45}%
  \let\sptunderscore\_\def\_{\sptunderscore\allowbreak}}
\setlength{\LTcapwidth}{\textwidth}
% A short table is kept on one page.  The build (code/scripts/keep_tables.py) ends the rows of such a table with \\*,
% which forbids a page break after a row, and sets \sptkeptpenalty to 10000 in front of the table.  This penalty is
% put here in the place of the breakpoint (penalty 0) that longtable leaves after every table, so that the paragraph
% that follows a kept table, its caption, begins on the same page; it is set back to 0 after the table.
\newcount\sptkeptpenalty
\sptkeptpenalty=0
\makeatletter
\patchcmd{\endlongtable}{\endgraf\penalty\z@\addvspace\LTpost}%
  {\endgraf\penalty\sptkeptpenalty\addvspace\LTpost\global\sptkeptpenalty=\z@}%
  {}{\PackageError{doc_header}{the end of longtable could not be patched}{The tables can still be typeset, but a short table may be separated from its caption.}}
\makeatother
% allow line breaks after slashes and hyphens inside table cells
\usepackage{xurl}
% code blocks: break over-long lines instead of running into the margin.
% Plain ``` blocks are typeset by pandoc as {verbatim}, highlighted ones
% (```bash) as {Highlighting}; both are re-declared as breakable fancyvrb
% environments.  The continuation line is marked with a small hooked arrow.
\usepackage{fvextra}
\RecustomVerbatimEnvironment{verbatim}{Verbatim}{breaklines,breakanywhere,fontsize=\small}
\makeatletter
\@ifundefined{Highlighting}{}{%
  \RecustomVerbatimEnvironment{Highlighting}{Verbatim}{breaklines,breakanywhere,fontsize=\small,commandchars=\\\{\}}}
\makeatother
% inline code (pandoc writes it as \texttt{...}): never hyphenate -- an inserted
% hyphen could be read as part of a file name -- but allow a line break after
% an underscore, slash, comma, colon or equals sign, so that long paths and
% column lists wrap inside the text width and inside table columns.
% Inside a formula (e.g. the name of a notebook symbol in an equation) the same
% command must first leave math mode, which \text of amsmath does.
\ExplSyntaxOn
\tl_new:N \l__spt_code_tl
\cs_new_protected:Npn \__spt_texttt:n #1
  {
    \tl_set:Nn \l__spt_code_tl {#1}
    \regex_replace_all:nnN { ([\/\,\:\=]) } { \1 \c{allowbreak} } \l__spt_code_tl
    \regex_replace_all:nnN { \c{_} } { \c{_} \c{allowbreak} } \l__spt_code_tl
    \group_begin: \ttfamily \hyphenchar \font = -1 \scan_stop: \l__spt_code_tl \group_end:
  }
\RenewDocumentCommand \texttt { m }
  {
    \mode_if_math:TF { \text { \__spt_texttt:n {#1} } } { \__spt_texttt:n {#1} }
  }
\ExplSyntaxOff
% A figure is written in the Markdown source as an image paragraph followed by
% a caption paragraph ("**Figure S1.** ..."), so that the caption is also
% visible on GitHub.  An image that starts a paragraph is centered.
\usepackage{graphicx}% pandoc loads it only for documents that contain an image
\NewCommandCopy\sptincludegraphics\includegraphics
\RenewDocumentCommand\includegraphics{O{} m}{%
  \ifvmode\begin{center}\sptincludegraphics[#1]{#2}\end{center}%
  \else\sptincludegraphics[#1]{#2}\fi}
% keep figures where they are referenced
\usepackage{float}
\floatplacement{figure}{H}
% running header with the document title
\usepackage{fancyhdr}
\pagestyle{fancy}
\fancyhf{}
\fancyhead[L]{\small PROJECT spacetimes, Revision\_SuperLoop}
\fancyhead[R]{\small\thepage}
\renewcommand{\headrulewidth}{0.4pt}
% 16x16 matrices: tighter columns so that a full matrix fits the text width
\setlength{\arraycolsep}{3.2pt}
% The documents of this project contain long tagged symbols, file paths and
% identifiers that cannot be hyphenated.  A paragraph is therefore allowed to
% have looser lines before TeX gives up and lets a line stick out of the margin.
\AtBeginDocument{\tolerance=2000 \emergencystretch=5em \hbadness=2000}
\ifLuaTeX
  \usepackage{selnolig}  % disable illegal ligatures
\fi
\IfFileExists{bookmark.sty}{\usepackage{bookmark}}{\usepackage{hyperref}}
\IfFileExists{xurl.sty}{\usepackage{xurl}}{} % add URL line breaks if available
\urlstyle{same}
\hypersetup{
  pdftitle={The stop of all running stages on 20 November 2026},
  pdfauthor={Prepared for Patrick Lee Nash, Ph.D.},
  colorlinks=true,
  linkcolor={blue},
  filecolor={Maroon},
  citecolor={blue},
  urlcolor={blue},
  pdfcreator={LaTeX via pandoc}}

\title{The stop of all running stages on 20 November 2026}
\usepackage{etoolbox}
\makeatletter
\providecommand{\subtitle}[1]{% add subtitle to \maketitle
  \apptocmd{\@title}{\par {\large #1 \par}}{}{}
}
\makeatother
\subtitle{Provenance document \texttt{P3} of PROJECT spacetimes,
Revision\_SuperLoop}
\author{Prepared for Patrick Lee Nash, Ph.D.}
\date{November 20, 2026}

\begin{document}
\maketitle

\hypertarget{what-this-document-is}{%
\section{1. What this document is}\label{what-this-document-is}}

On 20 November 2026 the author gave a new task. Its first instruction
was to stop all running stages and to create a provenance of the stop.
This document is that provenance. It records

\begin{itemize}
\tightlist
\item
  the two instructions of that day that concern the stop, verbatim
  (Section 3);
\item
  every stage of the work plan, what it was meant to do, and its state
  at the moment of the stop (Sections 4 and 5);
\item
  what was done to stop the stages, and what was found when the state
  was checked (Sections 6 and 8);
\item
  the one stage whose finished work was delivered after the stop with
  the author's permission, stage E, and how it was verified before
  delivery (Section 7);
\item
  what the stop does not change (Section 9), and the errors made during
  the stop (Section 10).
\end{itemize}

Every term used here is defined in Section 2. The document contains no
formula.

\hypertarget{the-terms-used-in-this-document}{%
\section{2. The terms used in this
document}\label{the-terms-used-in-this-document}}

\begin{longtable}[]{@{}
  >{\raggedright\arraybackslash}p{(\columnwidth - 2\tabcolsep) * \real{0.1651}}
  >{\raggedright\arraybackslash}p{(\columnwidth - 2\tabcolsep) * \real{0.8349}}@{}}
\toprule\noalign{}
\begin{minipage}[b]{\linewidth}\raggedright
Term
\end{minipage} & \begin{minipage}[b]{\linewidth}\raggedright
Meaning
\end{minipage} \\
\midrule\noalign{}
\endhead
\bottomrule\noalign{}
\endlastfoot
PROJECT spacetimes & the author's project. Its repository is
\texttt{https://github.com/once-ere/spacetimes.git}; its working folder
on the author's computer is
\texttt{/Users/nsh/Developer/github/spacetimes}, called \textbf{the Mac
folder} below. \\
Revision\_SuperLoop & the folder of the repository in which the stages
of Section 4 were carried out. \\
workspace & the cloud computer on which the programs and documents of
the stages were written and run. It holds a copy (a clone) of the
repository. \\
commit & a recorded state of the repository, named here by the first
seven hexadecimal digits of its identifier. \textbf{main} is the
principal line of commits. \textbf{origin/main} is the state of main on
the server of the repository, as last seen by the Mac folder. To
\textbf{push} is to send the commits of the Mac folder to the server;
the author does this himself. \\
bundle & a file that carries commits from the workspace to the Mac
folder. There they are added with a fast-forward merge: commits are
added, and nothing that is already there is changed or removed. \\
tracked file & a file whose content is recorded in the commits. \\
full check & the script \texttt{code/scripts/check\_all.sh} of
Revision\_SuperLoop. It rebuilds every program, data file, figure and
document of the folder from its sources and runs every check: 69 steps,
about 21 minutes. Its last line is \texttt{ALL\ PASS} when every step
succeeds. \\
program check & one of the numbered checks that a program of the folder
prints in its report (for example \texttt{CJ-12}), each with the word
PASS or FAIL. \\
stage & one part of the work plan in force before the stop; Section 4
lists them with their content. \\
task list & the list of tasks of the working session, which the author
sees beside the conversation. \\
watcher & a small Wolfram Language program that runs inside the author's
Wolfram application on the Mac. It executes requests that the workspace
places as files in the folder \texttt{.git-transfer/mma/} of the Mac
folder, and writes the answers there. It is a tool for running Wolfram
code, not a stage. \\
SFBS & the superfluid boson singlet state: the state of spin zero formed
by pairs of spinor charges. \\
CJT & the method of J. M. Cornwall, R. Jackiw and E. Tomboulis,
``Effective action for composite operators'', Phys. Rev.~D 10, 2428
(1974). \\
RK4, CVODE, IDA & RK4 is the classical Runge-Kutta method of fourth
order with a fixed step. CVODE and IDA are integrators with error
control of the SUNDIALS library, used here in the pure Rust version of
the author's repository and, for IDA, inside the Wolfram function
\texttt{NDSolve}. \\
Test 2 & the author's test problem of Revision\_SuperLoop: the
gauge-covariant spinor field on the 3+1 manifold and the SFBS, in eleven
items numbered {[}1{]} to {[}11{]}. \\
\end{longtable}

\hypertarget{the-instructions-of-20-november-2026}{%
\section{3. The instructions of 20 November
2026}\label{the-instructions-of-20-november-2026}}

At 17:21 (Pacific time) the author gave a new task. Its first
instruction reads, verbatim:

\begin{quote}
``First stop all running stages; create a provenance.''
\end{quote}

At 17:31 the author added, verbatim:

\begin{quote}
``it is OK to push stage E 's. completed work, if this does NOT
interfere with your current task''
\end{quote}

The rest of the new task is carried out in a new folder of the
repository and is not part of this record.

\hypertarget{the-stages-of-the-work-plan-and-their-state-at-the-stop}{%
\section{4. The stages of the work plan, and their state at the
stop}\label{the-stages-of-the-work-plan-and-their-state-at-the-stop}}

The work plan in force was that of 18 November 2026: stages A, B and C,
and the six stages D to I that the author added on 18 November at 20:17.
Table 1 lists each stage, what it was meant to do, its state at the
moment of the stop, and what the stop did to it.

\textbf{Table 1.} The stages of the work plan.

\begin{longtable}[]{@{}
  >{\raggedright\arraybackslash}p{(\columnwidth - 6\tabcolsep) * \real{0.0556}}
  >{\raggedright\arraybackslash}p{(\columnwidth - 6\tabcolsep) * \real{0.3333}}
  >{\raggedright\arraybackslash}p{(\columnwidth - 6\tabcolsep) * \real{0.4333}}
  >{\raggedright\arraybackslash}p{(\columnwidth - 6\tabcolsep) * \real{0.1778}}@{}}
\toprule\noalign{}
\begin{minipage}[b]{\linewidth}\raggedright
Stage
\end{minipage} & \begin{minipage}[b]{\linewidth}\raggedright
What it was meant to do
\end{minipage} & \begin{minipage}[b]{\linewidth}\raggedright
State at the stop
\end{minipage} & \begin{minipage}[b]{\linewidth}\raggedright
Action
\end{minipage} \\
\midrule\noalign{}
\endhead
\bottomrule\noalign{}
\endlastfoot
A & the notebook \texttt{3+1super-loop.nb} (at the top of the
repository) with every time integration done by IDA instead of RK4 &
complete; commit \texttt{8779db7} & none needed \\
B & the port of the programs of Test 2, items {[}2{]} to {[}11{]}, into
the same notebook & item {[}2{]} complete (commit \texttt{5dc688e});
items {[}3{]} to {[}11{]} not begun; a notebook chapter with the checks
named GS-17 to GS-20 had been written and set aside outside the
repository; it was never committed & stopped \\
C & the older programs that use integrators with a fixed step, redone
with CVODE in new folders & not begun & stopped \\
D & the creation of the SFBS from the spinor field by the method CJT &
complete; commit \texttt{32a4961}; pushed by the author (origin/main =
\texttt{32a4961} on the Mac at 17:30 on 20 November) & none needed \\
E & the effective potential of the SFBS in the method CJT (the analogue
of the ``Mexican hat'' potential of the complex Higgs doublet) & every
part written on 19 November: the program part (program checks CJ-12 to
CJ-15; the report of 19 November, 09:32, says ``15 of 15 checks
passed''), the new section 10 of the document of the stages D to I, the
figure, the data files, the symbol registry and the records of the
folder. \textbf{Not verified}: the full check had been started three
times; the first two runs failed through two errors of procedure of the
person doing the work, and the third run (started 19 November, 09:36:54)
stopped with the workspace after 17 of its 69 steps; its log ends at
09:41. \textbf{Not committed.} & verified and delivered after the
author's permission of 17:31 (Section 7) \\
F & the relative minima of this potential and the spontaneous breaking
of the symmetry & a draft of its program part (four program checks,
named CJ-16 to CJ-19) and a scratch calculation exist in the workspace
outside the repository; the draft was never compiled into the program &
stopped; the drafts are kept unchanged and are not continued \\
G & the expansion of the SFBS about the minimum (called SFBSv) & not
begun & stopped \\
H & the kinetic Lagrangian of the SFBS & not begun & stopped \\
I & the masses of the Goldstone bosons from SFBSv and the kinetic
Lagrangian, the diagonalization of the mass terms, the surviving
symmetry group, and its relation to U(1)\_em and the massless photon &
not begun & stopped \\
\end{longtable}

\hypertarget{stages-that-the-author-had-paused-earlier}{%
\section{5. Stages that the author had paused
earlier}\label{stages-that-the-author-had-paused-earlier}}

These were not running at the stop; they remain as the author left them.

\textbf{Table 2.} Stages paused or stopped earlier by the author.

\sptkeptpenalty=10000
\begin{longtable}[]{@{}
  >{\raggedright\arraybackslash}p{(\columnwidth - 4\tabcolsep) * \real{0.1712}}
  >{\raggedright\arraybackslash}p{(\columnwidth - 4\tabcolsep) * \real{0.5586}}
  >{\raggedright\arraybackslash}p{(\columnwidth - 4\tabcolsep) * \real{0.2703}}@{}}
\toprule\noalign{}
\begin{minipage}[b]{\linewidth}\raggedright
When
\end{minipage} & \begin{minipage}[b]{\linewidth}\raggedright
What
\end{minipage} & \begin{minipage}[b]{\linewidth}\raggedright
State
\end{minipage} \\
\midrule\noalign{}
\endhead
\bottomrule\noalign{}
\endlastfoot
12 November, about 15:20 & the main project (the 4+4 and 7+1
spacetimes): the independent verification pass, the correction of the
statement about fermions bound to a vortex line, a planned provenance,
the final check & paused by the author; not running since \\*
15 November, 12:36 & all stages of the main project then running &
paused by the author; not running since \\*
15 November, 13:25 & the folders \texttt{Revision\_2} and
\texttt{Revision\_2b} & stopped by the author's choice; they stay as
they are \\
\end{longtable}

\hypertarget{what-was-done-to-stop-the-stages}{%
\section{6. What was done to stop the
stages}\label{what-was-done-to-stop-the-stages}}

\begin{enumerate}
\def\labelenumi{\arabic{enumi}.}
\tightlist
\item
  \textbf{Processes in the workspace.} At 17:29 on 20 November no
  process of any stage ran in the workspace: it had been restarted, and
  the list of its processes contained none of the programs of the
  stages. There was nothing to end.
\item
  \textbf{Processes on the Mac.} No program of a stage runs on the Mac;
  the stages ran in the workspace. The watcher is not a stage and was
  not touched. It is needed for the new task, for which the author asked
  that the Wolfram application on his computer be used.
\item
  \textbf{The task list.} Every open task of the stages was marked
  ``STOPPED 2026-11-20''.
\item
  \textbf{The drafts.} The drafts of stage F and the set-aside notebook
  chapter of stage B were left where they were, unchanged, outside the
  repository. Nothing was deleted.
\item
  \textbf{No further work} is done on stages B, C, F, G, H and I, unless
  the author asks for it. The only work done after the stop on a stage
  of Table 1 is the verification and delivery of the finished work of
  stage E (Section 7), which the author permitted at 17:31 under the
  condition that it not interfere with the new task: it ran in the
  background while the new task began.
\end{enumerate}

\hypertarget{stage-e-verification-and-delivery}{%
\section{7. Stage E: verification and
delivery}\label{stage-e-verification-and-delivery}}

Stage E is the only stage of Table 1 on which work was done after the
stop, and that work was limited to verifying and delivering what had
been finished before the stop. Nothing of its content was changed.

\begin{enumerate}
\def\labelenumi{\arabic{enumi}.}
\tightlist
\item
  \textbf{The permission.} At 17:31 the author permitted the delivery of
  the finished work of stage E, if it did not interfere with the new
  task. It did not: the verification ran by itself in the background of
  the workspace while the new task began.
\item
  \textbf{The verification.} The full check of Revision\_SuperLoop was
  run once more from the start, on 20 November from 17:33:13 to
  17:52:42. Every one of its 69 steps succeeded; its last line is
  \texttt{ALL\ PASS}. The report of the program of stage E says ``15 of
  15 checks passed'' (the program checks CJ-1 to CJ-15; CJ-12 to CJ-15
  belong to stage E).
\item
  \textbf{The commit.} The files of stage E, as the full check rebuilt
  them from their sources, were committed unchanged as commit
  \texttt{0cefb61} (``Test 2, stage E: the effective potential of the
  superfluid boson singlet state by the method of Cornwall, Jackiw and
  Tomboulis''): 110 files.
\item
  \textbf{What stage E contains.} A new Section 10 of the document of
  the stages D to I of Test 2 (its file name: ``The superfluid boson
  singlet state by the method of Cornwall, Jackiw and Tomboulis'', .md,
  .tex, .pdf), with four theorems, two tables and one figure; the four
  program checks CJ-12 to CJ-15 of the program \texttt{t2\_cjt}; two
  data files and one figure; six entries of the symbol registry of the
  folder; and the records of the folder (its task record, symbol file,
  record of the calculations and README) brought up to date. The other
  files of the commit are those that the full check rebuilds every time
  (documents and figures whose content did not change).
\item
  \textbf{The transfer.} Commit \texttt{0cefb61} and the commit of this
  document are carried to the Mac folder together; the author pushes
  them.
\end{enumerate}

\hypertarget{the-state-of-the-repository-at-the-stop}{%
\section{8. The state of the repository at the
stop}\label{the-state-of-the-repository-at-the-stop}}

The repository on the Mac was examined at 17:30 on 20 November.

\begin{itemize}
\tightlist
\item
  main = origin/main = \texttt{32a4961}: the server and the Mac folder
  held the same commits.
\item
  No tracked file was changed or missing.
\item
  \texttt{git\ fsck}, the integrity check of git, found no missing and
  no damaged object.
\item
  It reported twelve leftover temporary files of git (34 MiB in total)
  in \texttt{.git/objects/pack/}, and a leftover lock file
  \texttt{.git/objects/maintenance.lock}. They were left by earlier git
  operations that could not finish removing their own temporary files,
  because the programs of the workspace are not allowed to delete files
  in the Mac folder. They do not affect the commits, the merges or the
  pushes. They were not removed.
\end{itemize}

These checks show that the content recorded in the repository is intact.
They do not judge the quality of that content.

After the stop, commit \texttt{0cefb61} (stage E) and the commit of this
document are carried to the Mac folder in one bundle and added there by
a fast-forward merge, which changes nothing that is already there; the
author then pushes them to the server.

\hypertarget{what-the-stop-does-not-change}{%
\section{9. What the stop does not
change}\label{what-the-stop-does-not-change}}

\begin{itemize}
\tightlist
\item
  Every commit up to \texttt{32a4961} stays as it is; nothing was
  rewritten or removed.
\item
  The folders of earlier revisions are not changed.
\item
  The new task of 20 November is carried out in a new folder of the
  repository, \texttt{Revision\_HS\_CJT\_and\_Kohn-Sham}. It changes
  nothing in Revision\_SuperLoop.
\end{itemize}

\hypertarget{errors-made-during-the-stop}{%
\section{10. Errors made during the
stop}\label{errors-made-during-the-stop}}

\begin{longtable}[]{@{}
  >{\raggedright\arraybackslash}p{(\columnwidth - 6\tabcolsep) * \real{0.0877}}
  >{\raggedright\arraybackslash}p{(\columnwidth - 6\tabcolsep) * \real{0.3947}}
  >{\raggedright\arraybackslash}p{(\columnwidth - 6\tabcolsep) * \real{0.1930}}
  >{\raggedright\arraybackslash}p{(\columnwidth - 6\tabcolsep) * \real{0.3246}}@{}}
\toprule\noalign{}
\begin{minipage}[b]{\linewidth}\raggedright
Error
\end{minipage} & \begin{minipage}[b]{\linewidth}\raggedright
What went wrong
\end{minipage} & \begin{minipage}[b]{\linewidth}\raggedright
How it was found
\end{minipage} & \begin{minipage}[b]{\linewidth}\raggedright
Correction
\end{minipage} \\
\midrule\noalign{}
\endhead
\bottomrule\noalign{}
\endlastfoot
E1 & The first draft of Table 2 of this document named another document
as the record of an earlier pause. The author's rule forbids sending the
reader to another document. & reading the draft against the author's
rules before the document was built & the reference was removed; Table 2
states the facts itself \\
\end{longtable}

\end{document}
